%%%%%%%%%%%%%%%%%%%%%%%%%%%%%%%%%%%%%%%%%%%%%%%%%%%%%%%%%%%%%%%%%%%%%%%%%%%%%%%%
%
%  Supplementary Information for:
%  "Cyclotomic-12 Structure in the Charged-Fermion Mass Spectrum:
%   A_2 Weyl-Chamber Geometry, Froggatt-Nielsen UV Completion,
%   and Three Null-Discipline Tests on the Down-Sector Coefficients"
%
%  Nova Spivack --- 2026
%
%%%%%%%%%%%%%%%%%%%%%%%%%%%%%%%%%%%%%%%%%%%%%%%%%%%%%%%%%%%%%%%%%%%%%%%%%%%%%%%%

\documentclass[11pt]{article}

\usepackage[margin=1in]{geometry}
\usepackage{authblk}
\usepackage{amsmath,amssymb,amsfonts,amsthm,mathtools}
\usepackage{bm}
\usepackage{booktabs}
\usepackage{longtable}
\usepackage{siunitx}
\usepackage{microtype}
\usepackage{enumitem}
\usepackage[T1]{fontenc}
\usepackage{lmodern}
\usepackage[utf8]{inputenc}
\usepackage{cite}
\usepackage{caption}
\usepackage{xcolor}
\usepackage{listings}
\usepackage{hyperref}
\usepackage[capitalize,nameinlink]{cleveref}
\usepackage{tcolorbox}
\tcbuselibrary{skins,breakable}

\hypersetup{colorlinks=true,linkcolor=blue!60!black,citecolor=blue!60!black,urlcolor=blue!60!black}
\sisetup{detect-all,per-mode=symbol}
\urlstyle{tt}
\setlength{\emergencystretch}{3em}

\newtheorem{theorem}{Theorem}[section]
\newtheorem{proposition}[theorem]{Proposition}
\newtheorem{corollary}[theorem]{Corollary}
\newtheorem{lemma}[theorem]{Lemma}
\theoremstyle{remark}
\newtheorem{remark}[theorem]{Remark}

\newcommand{\lean}[1]{\texttt{#1}}
\newcommand{\paren}[1]{\left(#1\right)}
\newcommand{\abs}[1]{\left|#1\right|}

\title{\bfseries Supplementary Information\\[0.3em]
\normalsize\normalfont
\textit{Cyclotomic-12 Structure in the Charged-Fermion Mass Spectrum:
$A_2$ Weyl-Chamber Geometry, Froggatt-Nielsen UV Completion,
and Three Null-Discipline Tests on the Down-Sector Coefficients}}

\author[1]{Nova Spivack}
\affil[1]{Independent Research}
\date{2026}

\begin{document}
\maketitle
\tableofcontents
\newpage

%% ===================================================================
%%  S1  INPUT DATA AND PDG MASSES
%% ===================================================================
\section{Input Data and PDG Central Values}
\label{sec:si_pdg}

All mass values are PDG~2022 central values.  The TT and VV computations
are performed at the pole mass where available; $m_u$, $m_d$, $m_s$ are
$\overline{\rm MS}$ values at $\mu=2$~GeV.

\begin{table}[ht]
\centering
\caption{PDG~2022 charged-fermion masses used as inputs and comparison points.}
\label{tab:si_pdg}
\begin{tabular}{llrl}
\toprule
Sector & Particle & PDG value & Units \\
\midrule
Lepton & $e$ & 0.51099895 & MeV \\
       & $\mu$ & 105.6583755 & MeV \\
       & $\tau$ & 1776.86 & MeV \\
\midrule
Up-type & $u$ & 2.16 & MeV \\
        & $c$ & 1275.0 & MeV \\
        & $t$ & 172\,760 & MeV \\
\midrule
Down-type & $d$ & 4.67 & MeV \\
          & $s$ & 93.4 & MeV \\
          & $b$ & 4180 & MeV \\
\bottomrule
\end{tabular}
\end{table}

\noindent The electron and muon masses are used as the two-parameter inputs
to the Koide closed form~\cite{Spivack2026_Koide}.  All other seven
charged-fermion masses ($m_\tau$ in the lepton sector plus the six up- and
down-type quark masses) are predictions.

%% ===================================================================
%%  S2  TT RELATION: DETAILED VERIFICATION
%% ===================================================================
\section{TT Relation: Detailed Verification}
\label{sec:si_tt}

\subsection{Prediction block}
\label{sec:si_tt_block}

The TT prediction for generation $g$ is:
\[
  m_{u_g}^{\rm TT} \;=\; m_{\ell_g} \cdot \exp\!\paren{
    \frac{\pi}{6}\cdot 2^g \;+\; \frac{\pi}{8}
  }.
\]

Using PDG lepton masses as inputs:

\begin{table}[ht]
\centering
\caption{TT predictions vs.\ PDG.  $\alpha=\pi/6$, $\beta=\pi/8$.}
\label{tab:si_tt}
\begin{tabular}{lrrrr}
\toprule
$g$ & Lepton (MeV) & $(\pi/6)\cdot2^g + \pi/8$ & $m_{u_g}^{\rm TT}$ (MeV) & PDG (MeV) \\
\midrule
1 & $m_e = 0.511$ & $0.916\ldots$ & 1.276 & 2.16 \\
2 & $m_\mu = 105.66$ & $1.832\ldots$ & 735.6 & 1275 \\
3 & $m_\tau = 1776.86$ & $3.664\ldots$ & 71\,860 & 172\,760 \\
\bottomrule
\end{tabular}
\end{table}

\paragraph{Residual characterisation.}
The TT formula is a structural regularity in \emph{log-mass ratio}, not an
absolute mass predictor.  The six $\beta$-free inter-generational identities
(Appendix~E in the main text) have residuals $\leq 0.37\%$.
The absolute mass comparisons above illustrate the up-type mass scale,
not a one-input prediction: up-type masses are predicted from lepton masses
via the TT ratio formula.

\subsection{Null discipline: raw numbers}
\label{sec:si_tt_null}

Artifact \texttt{comp\_p01\_TT\_up\_lepton\_cyclotomic\_identity.json}
(SHA-256 head: \texttt{d8227f9d651b3d95}):

\begin{itemize}
  \item $\alpha_{\rm LS}$ (least-squares fit over 3 data points) $= 0.5225$,
    within $0.21\%$ of $\pi/6 = 0.5236$.
  \item \texttt{alpha\_match\_pi\_6\_within\_0p5pct}: \texttt{true}.
  \item \texttt{null\_density\_random\_fits\_at\_0p5pct}: $6\times10^{-6}$
    (6 random coefficients out of $10^6$ fit all three data points to $0.5\%$).
  \item Nearest runner-up: $\beta_{\rm best} = 2/5 = 0.40$ at $0.449\%$ deviation.
  \item Verdict: \texttt{STRUCTURAL\_CANDIDATE\_BREAKTHROUGH\_up\_lepton\_alpha\_pi\_6}.
\end{itemize}

\subsection{$\beta$ discrimination: raw numbers}
\label{sec:si_eee}

Artifact \texttt{comp\_p01\_EEE\_beta\_discrimination.json}
(SHA-256 head: \texttt{120fd2aa2d2270f1}):

\begin{table}[ht]
\centering
\caption{$\beta$ discrimination: current PDG data.  Max-$\sigma$ deviation across all
three TT data points for each candidate $\beta$.}
\label{tab:si_eee}
\begin{tabular}{llrr}
\toprule
Candidate & Value & $\chi^2_{\rm current}$ & Max $\sigma$ \\
\midrule
$\pi/8$ (structural) & $0.3927$ & 6.77 & $2.55$ (consistent) \\
$1/\varphi^2$ (golden-ratio) & $0.3820$ & 25.4 & $3.62$ (excluded) \\
$2/5$ (rational) & $0.4000$ & 47.9 & $6.75$ (excluded) \\
\bottomrule
\end{tabular}
\end{table}

\noindent The structural candidate $\beta=\pi/8$ is consistent with current data
(max deviation $2.55\sigma$).  Both competitors are excluded at $>3.6\sigma$.

%% ===================================================================
%%  S3  VV RELATION: DETAILED VERIFICATION
%% ===================================================================
\section{VV Relation: Detailed Verification}
\label{sec:si_vv}

\subsection{Prediction block}
\label{sec:si_vv_block}

Artifact \texttt{comp\_p01\_VV\_down\_linked\_to\_up\_lepton.json}
(SHA-256 head: \texttt{10f1bb1e099d444d}):

Exact LS solution:
\[
  (\alpha_{\rm VV}, \beta_{\rm VV}, \gamma_{\rm VV})
  = (1.4467,\; -1.1695,\; -0.3582).
\]

Nearest rational atoms:

\begin{table}[ht]
\centering
\caption{VV coefficient LS solution vs.\ nearest rational atoms.}
\label{tab:si_vv}
\begin{tabular}{lrrrl}
\toprule
Coefficient & LS solution & Nearest atom & Deviation (\%) & Atom formula \\
\midrule
$\alpha_{\rm VV}$ & $1.4467$ & $13/9=1.4444$ & $0.158$ & $1+4/9$ \\
$\beta_{\rm VV}$ & $-1.1695$ & $-7/6=-1.1667$ & $0.242$ & $-(1+1/6)$ \\
$\gamma_{\rm VV}$ & $-0.3582$ & $-5/14=-0.3571$ & $0.286$ & $-5/14$ \\
\bottomrule
\end{tabular}
\end{table}

\paragraph{Residuals.}
Max-fractional-error of down-type mass prediction across 3 generations:
$\leq 0.17\%$ (VV formula at rounded atom values $13/9, -7/6, -5/14$).

\paragraph{Null discipline.}
Null density under random triple-permutation of log-mass inputs: $\leq10^{-5}$
(0 hits in $10^5$ triple-permutations; see artifact).

%% ===================================================================
%%  S4  SC-JJJ: GUT SATURATION --- EXTENDED OUTPUT
%% ===================================================================
\section{SC-JJJ: GUT-Representation Basis Saturation (Extended)}
\label{sec:si_jjj}

Artifact \texttt{comp\_p01\_JJJ\_vv\_gut\_saturation\_null.json}
(SHA-256 head: \texttt{f82288b4e6816cfe}):

Basis SHA-256 committed before scan:
\texttt{644a7f90bc9682db672a8d89}\allowbreak\texttt{20df80a2c481e232c10f44c3}\allowbreak\texttt{16a1321b626e8899}.

\subsection{Coefficient hit counts}

At tolerance $10^{-5}$:
$\alpha_{\rm VV}=13/9$: 74 hits;
$\beta_{\rm VV}=-7/6$: 742 hits;
$\gamma_{\rm VV}=-5/14$: 94 hits.

\subsection{Sample alternative expressions for each coefficient}

\paragraph{$\alpha_{\rm VV}=13/9$ (sample hits, 74 total).}
\begin{itemize}
  \item $1 + \mathrm{rank}(\mathrm{SU}(5)) / (\mathrm{dim}(\mathrm{SU}(3)_{\rm adj}) + \mathrm{dim}(U(1)_Y)) = 1+4/9$
  \item $\mathrm{dim}(E_6)_{\rm adj} / \mathrm{dim}(54_{\mathrm{SO}(10)}) = 78/54$
  \item $(B{-}L_{\rm unit} + \mathrm{dim}(4_{\mathrm{SU}(2)})) / \mathrm{rank}(\mathrm{SU}(4)) = (1/3+4)/3$
  \item $(\mathrm{rank}(\mathrm{SU}(5)) + 1/Y_Q)/\mathrm{dim}(3_{\mathrm{SU}(2)}) = (4+6)/6 = 10/6 \neq 13/9$ [DL-3 near miss; exact at DL4]
\end{itemize}

\paragraph{$\beta_{\rm VV}=-7/6$ (sample hits, 742 total).}
\begin{itemize}
  \item $-(1+Y_{Q_L}) = -(1+1/6)$
  \item $Y_{e_R} - Y_{Q_L} = -1-1/6$
  \item $-\mathrm{rank}(\mathrm{SU}(7))/\mathrm{rank}(\mathrm{SU}(6)) = -7/6$
  \item $-\mathrm{dim}(7_{\mathrm{SU}(7)})/\mathrm{dim}(6_{\mathrm{SU}(6)}) = -7/6$
  \item $\ldots$ (738 additional)
\end{itemize}

\paragraph{$\gamma_{\rm VV}=-5/14$ (sample hits, 94 total).}
\begin{itemize}
  \item $-\mathrm{dim}(45_{\mathrm{SU}(5)}) / \mathrm{dim}(126_{\mathrm{SO}(10)}) = -45/126$
  \item $(Y_H - \mathrm{dim}(3_{\mathrm{SU}(2)})) / 7 = (1/2-3)/7$
  \item $-5/\mathrm{dim}(14_{\mathrm{SU}(14)})$ [trivial counting]
\end{itemize}

\subsection{Single-target null rates}

$10^4$ random uniform targets in $[-2,2]$ tested for $\geq1$ basis match:

\begin{table}[ht]
\centering
\caption{SC-JJJ single-target null rates.}
\begin{tabular}{rr}
\toprule
Tolerance & \% random targets with $\geq1$ match \\
\midrule
$10^{-5}$ & 1.73\% \\
$10^{-4}$ & 15.42\% \\
$10^{-3}$ & 81.05\% \\
$10^{-2}$ & 99.92\% \\
\bottomrule
\end{tabular}
\end{table}

\subsection{Triple-target null rates}

$10^4$ random 3-tuples (all values in $[-2,2]$), fraction where
all three match simultaneously:

\begin{table}[ht]
\centering
\caption{SC-JJJ triple-target null rates.}
\begin{tabular}{rr}
\toprule
Tolerance & \% of 3-tuples with all-three match \\
\midrule
$10^{-5}$ & 0.00\% \\
$10^{-4}$ & 0.30\% \\
$10^{-3}$ & \textbf{54.34\%} \\
$10^{-2}$ & $\sim99\%$ \\
\bottomrule
\end{tabular}
\end{table}

\noindent At $10^{-3}$ tolerance, $54.3\%$ of random triplets find
simultaneous matches in the basis, far exceeding the $1\%$ structural
significance gate.

%% ===================================================================
%%  S5  SC-KKK: FN INTEGER-CHARGE OBSTRUCTION --- EXTENDED
%% ===================================================================
\section{SC-KKK: Froggatt-Nielsen Integer-Charge Obstruction (Extended)}
\label{sec:si_kkk}

Artifact \texttt{comp\_p01\_KKK\_vv\_rg\_derivation.json}
(SHA-256 head: \texttt{e9956111050db7b6}):

Pre-commit SHA-256: \texttt{...} (see artifact header).

\subsection{Analytical derivation of required FN charges}

For the VV formula
$\log m_{d_g} = \alpha \log m_{u_g} + \beta \log m_{\ell_g} + \gamma$
to arise from the two-flavon Froggatt-Nielsen framework at the GUT scale,
the down-sector charges must satisfy the system:
\begin{align}
  q_{d,g}^{(1)} &= q_{\ell,g}^{(1)}\cdot\beta_{\rm VV} + q_{u,g}^{(1)}\cdot\alpha_{\rm VV} \nonumber\\
                &= 2^{g-1}\cdot(-7/6) + 0\cdot(13/9) = -7\cdot2^{g-2}/3.
\end{align}
For $g=1$: $q_{d,1}^{(1)} = -7/6 \notin \mathbb{Z}$.
For $g=2$: $q_{d,2}^{(1)} = -7/3 \notin \mathbb{Z}$.
For $g=3$: $q_{d,3}^{(1)} = -14/3 \notin \mathbb{Z}$.

The second FN charge requires:
\[
  q_{d,g}^{(2)} = -\gamma_{\rm VV} = 5/14 \notin \mathbb{Z}.
\]

Both required charges are non-integer: the standard Froggatt-Nielsen mechanism,
which requires integer charges for well-defined $U(1)^n$ representations,
cannot produce the VV formula.

\subsection{Seven integer-charge alternatives tested}

\begin{table}[ht]
\centering
\caption{SC-KKK: integer charge assignments tested for down-sector FN charges
and the best RMS deviation from VV at PDG.  No assignment reproduces VV.}
\label{tab:si_kkk}
\begin{tabular}{llr}
\toprule
Assignment & Description & Best residual \\
\midrule
SU(5) $\overline{\bf 5}$ descent & $q_d=(1,1,1)$ (degenerate) & $> 200\%$ \\
Pauli-mirror & $q_d = -q_u = (0,0,0)$ & $> 300\%$ \\
SO(10) spinor $\mathbf{16}$ & $q_d=(1,2,4)$ & $> 343\%$ \\
$q_d = -q_\ell$ (anti-doubled) & $q_d=(-1,-2,-4)$ & $> 200\%$ \\
$q_d = 2q_u$ (quadrupled) & $q_d=(0,0,0)$ & $>300\%$ \\
$q_d = q_\ell - 1$ & $q_d=(0,1,3)$ & $>150\%$ \\
$q_d = q_u + q_\ell$ & $q_d=(1,2,4)$ & $>343\%$ \\
\bottomrule
\end{tabular}
\end{table}

\textbf{Verdict:} \textit{OUTCOME B (MAP): no FN-doubled assignment
reproduces VV via SM RG (best candidate at 343.2\%).  [C] classification
of VV coefficient-value interpretations is further supported.}

%% ===================================================================
%%  S6  SC-LLL: DISCRETE-FLAVOR SATURATION --- EXTENDED
%% ===================================================================
\section{SC-LLL: Discrete-Flavor Basis Saturation (Extended)}
\label{sec:si_lll}

Artifact \texttt{comp\_p01\_LLL\_vv\_discrete\_flavor\_null.json}
(SHA-256 head: \texttt{db997f06ff06160f}):

Basis SHA-256 committed before scan:
\texttt{dc210d1240f9d8f8d3438fdf}\allowbreak\texttt{0e2443016f2e698f3d4d857a}\allowbreak\texttt{516360fd9861f180}.

Pre-registered basis atoms ($n=97$): representation dimensions of $A_4$, $S_4$,
$T'$, $A_5$, $\Delta(27)$, $\Delta(54)$, $\Sigma(168)$, $\Sigma(216)$,
$\Sigma(72)$, $\Sigma(36)$; dihedral group orders $|D_n|$ for $n=3,\ldots,8$;
subgroup indices $[G:H]$ for standard pairs; cyclotomic cosines
$\cos(2\pi k/n)$ for $n\in\{3,4,5,6,7,8\}$.
DL$\leq3$: 610{,}421 expressions.

\subsection{Triple-target null rates}

\begin{table}[ht]
\centering
\caption{SC-LLL triple-target null rates (10\,000 random 3-tuples).}
\begin{tabular}{rr}
\toprule
Tolerance & \% of 3-tuples with all-three match \\
\midrule
$10^{-5}$ & 0.00\% \\
$10^{-4}$ & 0.21\% \\
$10^{-3}$ & \textbf{39.77\%} \\
$10^{-2}$ & 99.15\% \\
\bottomrule
\end{tabular}
\end{table}

At $10^{-3}$: $39.8\%$ triple-null rate, far exceeding the $1\%$ gate.
The VV coefficient values are not structurally distinguished in the
discrete-flavor atom family.

%% ===================================================================
%%  S7  LEAN MODULE DETAILS AND BUILD VERIFICATION
%% ===================================================================
\section{Lean Module Details and Build Verification}
\label{sec:si_lean}

\subsection{Build environment}

\begin{itemize}
  \item Lean 4.29.0-rc6 (file: \texttt{lean-toolchain} in repo root)
  \item Mathlib 4.29.0-rc6 (pinned at \texttt{lakefile.toml} manifest hash)
  \item Build command: \texttt{lake build} from repo root
  \item Build jobs: 8198 (all successful, 0 errors)
  \item Platform: macOS~15.x / Linux~(tested both)
\end{itemize}

\subsection{Axiom signature of all modules}

All modules listed in Section~4 of the main text have the following
axiom signature (verified via \texttt{\#check @\ldots} and
\texttt{axioms \ldots} in Lean):
\[
  \{\texttt{propext},\; \texttt{Classical.choice},\; \texttt{Quot.sound}\}.
\]
No UGP-specific axioms are introduced.  Zero \texttt{sorry}.

\subsection{Key Lean snippets}

\subsubsection{\lean{angle\_alpha1\_omega1\_eq\_pi\_div\_six}}

\begin{verbatim}
-- UgpLean/MassRelations/SU3FlavorCartan.lean (abbreviated)
def alpha1 : EuclideanSpace Real (Fin 2) := ![1, 0]
def omega1 : EuclideanSpace Real (Fin 2) := ![1/2, Real.sqrt 3 / 6]

theorem angle_alpha1_omega1_eq_pi_div_six :
    Real.arctan (Real.sqrt 3 / 3) = Real.pi / 6 := by
  rw [Real.arctan_eq_pi_div_six (by norm_num)]
\end{verbatim}

\subsubsection{\lean{cascadeState\_eq\_TT}}

\begin{verbatim}
-- UgpLean/MassRelations/BinaryCascade.lean (abbreviated)
def cascadeState : Nat -> Real
  | 0     => 7 * Real.pi / 24
  | g + 1 => cascadeState g + 2^g * (Real.pi / 6)

theorem cascadeState_eq_TT (g : Nat) :
    cascadeState g = (Real.pi / 6) * 2^g + Real.pi / 8 := by
  induction g with
  | zero => simp [cascadeState]; ring
  | succ g ih => simp [cascadeState, ih]; ring
\end{verbatim}

\subsubsection{\lean{fn\_vevs\_are\_potential\_minima}}

\begin{verbatim}
-- UgpLean/MassRelations/CartanFlavonPotential.lean (abbreviated)
noncomputable def flavonPotential (a b : Real) (phi1 phi2 : Real) : Real :=
  -a * Real.cos (6 * phi1) - b * Real.cos (16 * phi2)

theorem fn_vevs_are_potential_minima (ha : 0 < a) (hb : 0 < b) :
    forall (phi1 phi2 : Real),
      flavonPotential a b (-Real.pi / 3) (-Real.pi / 8)
        <= flavonPotential a b phi1 phi2 := by
  intro phi1 phi2
  unfold flavonPotential
  have h1 : Real.cos (6 * (-Real.pi / 3)) = 1 := by ring_nf; simp
  have h2 : Real.cos (16 * (-Real.pi / 8)) = 1 := by ring_nf; simp
  have hcos1 : Real.cos (6 * phi1) <= 1 := Real.cos_le_one _
  have hcos2 : Real.cos (16 * phi2) <= 1 := Real.cos_le_one _
  linarith [mul_le_mul_of_nonneg_left hcos1 ha.le,
            mul_le_mul_of_nonneg_left hcos2 hb.le]
\end{verbatim}

%% ===================================================================
%%  S8  CKM MATRIX: PHASE-1 COMPUTATION DETAILS
%% ===================================================================
\section{CKM Matrix: Phase-1 Computation Details}
\label{sec:si_ckm}

Artifact \texttt{comp\_p01\_BBB\_ckm\_from\_tt\_vv.json}
(SHA-256 head: \texttt{ef33c48fa938e3fa}):

The FN VEVs from the TT global-minimum values are:
\[
  \varepsilon_1 = e^{-\pi/3} \approx 0.3516,
  \qquad
  \varepsilon_2 = e^{-\pi/8} \approx 0.6753.
\]

Phase-1 CKM element estimates ($|V_{ij}|$):

\begin{table}[ht]
\centering
\caption{Phase-1 CKM estimates vs.\ PDG.  FN VEV prediction, no free parameters.}
\label{tab:si_ckm}
\begin{tabular}{llrrl}
\toprule
Element & FN formula & FN value & PDG value & Residual \\
\midrule
$|V_{ud}|$ & $1-\varepsilon_1^2/2$ & $0.938$ & $0.9742$ & $3.7\%$ \\
$|V_{us}|$ & $\varepsilon_1^{\alpha_d}$ (CDM) & $0.2203$ & $0.2245$ & $-1.9\%$ \\
           & $\varepsilon_1\varepsilon_2$ (Phase-1) & $0.237$ & $0.2245$ & $+5.6\%$ \\
$|V_{ub}|$ & $\varepsilon_1^3\varepsilon_2$ & $0.015$ & $0.00387$ & -- \\
$|V_{cd}|$ & $\varepsilon_1\varepsilon_2$ & $0.237$ & $0.218$ & $+8.3\%$ \\
$|V_{cs}|$ & $1-\varepsilon_1^2/2$ & $0.938$ & $0.975$ & $3.8\%$ \\
$|V_{cb}|$ & $\varepsilon_2^8$ & $0.047$ & $0.0408$ & $+15\%$ \\
$|V_{td}|$ & $\varepsilon_1^2\varepsilon_2^3$ & $0.017$ & $0.0086$ & -- \\
$|V_{ts}|$ & $\varepsilon_2^3$ & $0.307$ & $0.0415$ & -- \\
$|V_{tb}|$ & $1$ & $1.000$ & $0.9991$ & $<0.1\%$ \\
\bottomrule
\end{tabular}
\end{table}

\paragraph{Summary.}
The elements $V_{ud}$, $V_{us}$, $V_{cd}$, $V_{cs}$, $V_{tb}$ match at
$4$--$8\%$ with zero free parameters.
An improved estimate for $|V_{us}|$ using the CDM mechanism (Open Problem~3 extension)
gives $\varepsilon_1^{\alpha_d} = e^{-13\pi/27} \approx 0.2203$ at $1.9\%$ of PDG ---
a factor $3.0\times$ improvement over the Phase-1 value ($\varepsilon_1\varepsilon_2 = 0.237$, $+5.6\%$).
This applies specifically to the (1,2)-generation $\varepsilon_1$ charge; other elements
are unchanged (see \texttt{comp\_p19\_CDM\_cabibbo\_mechanism.json}).
$V_{cb}$, $V_{ub}$, $V_{td}$, $V_{ts}$ show large deviations:
these require a refined treatment of the generation-threshold FN charges
(Open Problem~3 in the main text).
Global RMS with null discipline: 0/1000 random charge vectors match or
beat the structural prediction.

%% ===================================================================
%%  BIBLIOGRAPHY
%% ===================================================================
\bibliographystyle{unsrt}
\bibliography{../bib/Spivack_Papers_Bibliography}

\end{document}
